% Supported v0.5 teaching-statement example.
\documentclass[type=teaching]{careerdossier-statement}
% examples/profiles/profile-academic.tex
%
% Shared profile metadata used by the academic-CV example and the corresponding
% industry-resume smoke fixture. Keeping it independent of document layout
% proves that both document families consume the same \CDossierSetup interface.

\CDossierSetup{
  name     = {Ada Lovelace},
  headline = {Researcher in Analytical Computing},
  email    = {ada.lovelace@example.com},
  location = {Toronto, Ontario, Canada},
  website  = {example.com/ada-lovelace},
  %% The bare Google Scholar identifier, expanded to
  %% scholar.google.com/citations?user=ada-example for both the displayed text
  %% and the link. Pasting the whole address works too and is left as written.
  %% See docs/API.md, "Accepted forms for the web-profile keys".
  scholar  = {ada-example},
  orcid    = {0000-0002-1825-0097},
  affiliation = {Example University},
}

\CDossierStatementSetup{
  subtitle={Teaching transparent reasoning through computational practice},
  application-context={Faculty application, Department of Computational Science}
}
\begin{document}
\MakeCDossierStatementHeader

\section*{Teaching Purpose and Commitments}

I teach students to make their reasoning visible: to identify assumptions,
connect formal ideas to evidence, and explain how they know a result is
trustworthy. In computational courses, this means treating code, mathematical
argument, and communication as parts of one intellectual practice. My goal is
not simply that students can reproduce a technique, but that they can decide
when it applies, test its limits, and communicate their decisions to others.

Students enter a course with different disciplinary backgrounds and different
levels of confidence. I therefore organize learning around recurring questions
rather than presumed familiarity with a particular tool. What is the claim?
What evidence would change our conclusion? Which approximation is being made?
This shared language allows experienced programmers and first-time coders to
contribute to the same discussion without lowering the conceptual standard.

\section*{Teaching in Practice}

I use worked examples as objects of inquiry rather than scripts to copy. After
introducing a numerical method, I ask students to predict how it will behave on
a deliberately difficult case, compare that prediction with an experiment, and
explain any mismatch. Small groups record assumptions and testing strategies
before they run code. The class then compares approaches, including apparently
successful results that are supported by weak reasoning.

In project work, students move through short cycles of proposal, implementation,
peer review, and revision. A project brief defines the scientific question and
quality criteria while leaving meaningful choices open. Students review one
another's documentation and tests, which makes audience and reproducibility
concrete. I provide checkpoints early enough that feedback can change the work,
not merely justify a final grade.

I also connect technical content to authentic decisions. A lesson on error
estimation may use the question of whether two simulations support the same
conclusion; a lesson on data cleaning may examine how exclusions alter a
reported pattern. These contexts help students see that methodological choices
carry consequences beyond producing a correct-looking answer.

\section*{Assessment, Inclusion, and Responsiveness}

Assessment should reveal how students reason and give them opportunities to act
on feedback. I combine low-stakes retrieval, annotated problem solving, code
reviews, and larger integrative projects. Rubrics distinguish conceptual
reasoning, implementation quality, validation, and communication so that a
single error does not conceal what a student understands. Short reflections ask
students to identify the strongest evidence for their conclusion and the next
test they would run.

Inclusive teaching begins with transparent expectations and multiple routes
into participation. I provide examples of successful work, explain the purpose
of each task, and make key materials available in accessible formats before
class. Students may prepare questions in writing, discuss them in pairs, or
contribute to a whole-class exchange. Structured roles in group work rotate so
that technical authority does not settle around the most confident speaker.

I use student work, anonymous check-ins, and assessment patterns to revise my
teaching. When several students make the same error, I treat it as evidence
about the design of an explanation or task. I then reteach the idea through a
different representation and document the change for the next course offering.

\section*{Contribution to Example University}

At Example University, I would contribute to courses in numerical analysis,
scientific computing, and research methods, while developing project pathways
that connect undergraduate and graduate learners with open research. I am
especially interested in collaborating on shared assessment resources and in
mentoring students who want to bridge mathematical foundations with durable
software practice. My aim is a learning environment in which rigor is explicit,
feedback is usable, and every student can become a more independent reasoner.

\end{document}
